\subsection{Summary \quad / \quad \textthai{บทสรุป}}
\begin{paracol}{2}
In this chapter, we have navigated the transition from classical spectral analysis to data-driven deep learning for signal detection. We established that:
\begin{itemize}
    \item \textbf{1D Convolutions} serve as adaptive, learned cross-correlation filters that outperform manual template matching in non-stationary noise.
    \item \textbf{STFT and Spectrograms} provide a bridge between time and frequency, but deep models can natively process raw strain data to bypass the Heisenberg resolution trade-off.
    \item \textbf{SNR Robustness}: DL models achieve high accuracy at noise floors where traditional methods yield excessive false positives.
\end{itemize}

\switchcolumn

\begin{thai}
ในบทนี้ เราได้สำรวจการเปลี่ยนผ่านจากการวิเคราะห์สเปกตรัมแบบคลาสสิกไปสู่การเรียนรู้เชิงลึกที่ขับเคลื่อนด้วยข้อมูลสำหรับการตรวจจับสัญญาณ โดยมีข้อสรุปดังนี้:
\begin{itemize}
    \item \textbf{คอนโวลูชัน 1 มิติ} ทำหน้าที่เป็นตัวกรองสหสัมพันธ์แบบไขว้ที่ปรับตัวได้และผ่านการเรียนรู้มาแล้ว ซึ่งทำงานได้ดีกว่าการจับคู่เทมเพลตด้วยตนเองในสภาวะที่มีสัญญาณรบกวนไม่คงที่
    \item \textbf{STFT และสเปกโทรแกรม} เป็นสะพานเชื่อมระหว่างเวลาและความถี่ แต่โมเดลเชิงลึกสามารถประมวลผลข้อมูลดิบได้โดยตรงเพื่อก้าวข้ามข้อจำกัดการแลกเปลี่ยนความละเอียดของไฮเซนเบิร์ก
    \item \textbf{ความทนทานต่อ SNR}: โมเดล DL บรรลุความแม่นยำสูงในระดับสัญญาณรบกวนพื้นฐานที่วิธีการดั้งเดิมมักให้ผลลัพธ์ที่เป็น False Positive จำนวนมาก
\end{itemize}
\end{thai}
\end{paracol}
